\section{Complex optical constants: attenuation and phase}
For a monochromatic plane wave travelling into a homogeneous material,
\begin{align}
E(z,t)&=\Rea\!\left\{E_0\ee^{\ii k_0\bn z}\ee^{-\ii\omega t}\right\},
&\bn&=n+\ii\kappa,\\
\widetilde E(z)&=E_0\underbrace{\ee^{-k_0\kappa z}}_{\text{amplitude decay}}
\underbrace{\ee^{\ii k_0nz}}_{\text{phase accumulation}}.
\end{align}
The intensity of one forward wave in the same homogeneous medium is proportional to $|\widetilde E|^2$, so
\begin{equation}
\boxed{\frac{I(z)}{I(0)}=\ee^{-\alpha z},\qquad
\alpha=2k_0\kappa=\frac{4\pi\kappa}{\lambda_0},\qquad
\ell_I=\frac{1}{\alpha}.}
\end{equation}
The amplitude decay length is $1/(k_0\kappa)=2\ell_I$. These lengths must not be confused.

Relative to an equal geometrical path $d$ in vacuum, the material changes the phase by
\begin{equation}
\Delta\phi_{\rm mat}=k_0(n-1)d.
\end{equation}
In EUV notation, $\bn=1-\delta+\ii\beta$:
\begin{equation}
\frac{\widetilde E_{\rm material}(d)}{\widetilde E_{\rm vacuum}(d)}
=\ee^{-k_0\beta d}\ee^{-\ii k_0\delta d}.
\end{equation}
Thus an absorber is a complex optical element: it attenuates and changes phase. Even when transmission is weak, its interfaces and the underlying reflector affect the reflected complex amplitude.

% \point{A binary black-paint model retains only a real amplitude mask. A local reflective model instead has $r(x,y)=|r(x,y)|\ee^{\ii\phi_r(x,y)}$. Where a scalar local model is adequate, a phase change gives an equivalent wavefront error $\Delta W=\Delta\phi_r/k_0$. At a reflection zero the phase is undefined; an unwrapped phase and a reference are needed elsewhere.}

% The propagation factor alone does not include entry/exit reflections. For oblique propagation through a planar film, use $\ee^{\ii k_zd}$ rather than replacing $z$ by a guessed ray-path length. The conserved transverse wave number and boundary phases are part of the thin-film problem. For a patterned EUV absorber, diffraction couples different transverse components; a stack calculation is only the planar building block.

% \textbf{Illustrative calculation, not a material measurement.} If $\lambda_0=13.5\,\mathrm{nm}$, $\delta=0.05$, $\beta=0.02$, and $d=20\,\mathrm{nm}$, then $I(d)/I(0)=\ee^{-4\pi\beta d/\lambda_0}\simeq0.689$ and $\Delta\phi_{\rm mat}\simeq-0.465\,\mathrm{rad}$. These chosen constants demonstrate the two simultaneous effects; they are not assigned to Mo, Si, or a particular absorber.

\section{Maxwell boundary conditions and single-interface amplitudes}
\subsection{Back to Griffiths: Tangential fields are continuous}
From Griffiths we know,
\begin{equation}
\nabla\times\mathbf E=\ii\omega\mathbf B,
\quad \nabla\times\mathbf H=\mathbf J_f-\ii\omega\mathbf D,
\quad \nabla\cdot\mathbf D=\rho_f,
\quad \nabla\cdot\mathbf B=0.
\end{equation}
Take a rectangular loop of length $L$ tangential to the interface and height $\varepsilon$ normal to it. Faraday's integral law has an electric circulation $(E_{1,\tau}-E_{0,\tau})L+O(\varepsilon)$ and a flux term $O(L\varepsilon)$ for bounded fields. Letting $\varepsilon\to0$ gives equality of the tangential electric fields. Ampere--Maxwell gives the corresponding magnetic relation, except that an imposed surface-current sheet can remain finite. With $\widehat{\mathbf n}$ directed from medium 0 to medium 1,
\begin{equation}
\boxed{\widehat{\mathbf n}\times(\mathbf E_1-\mathbf E_0)=0,
\qquad \widehat{\mathbf n}\times(\mathbf H_1-\mathbf H_0)=\mathbf K_f.}
\label{eq:bc}
\end{equation}
We take $\mathbf K_f=0$. A vanishing pillbox also gives
$\widehat{\mathbf n}\cdot(\mathbf D_1-\mathbf D_0)=\sigma_f$ and
$\widehat{\mathbf n}\cdot(\mathbf B_1-\mathbf B_0)=0$.
Normal $E$ need not be continuous, because the permittivity can jump. These are conditions on \emph{total fields} on each side, not on the incident wave alone.

Using a common fixed tangential electric axis, at the boundary
\begin{equation}
E_i+E_r=E_t,\qquad r=\frac{E_r}{E_i},\quad t=\frac{E_t}{E_i}
\quad\Longrightarrow\quad\boxed{1+r=t.}
\label{eq:Ebc}
\end{equation}
% This is a field-amplitude equation, not an equation for optical powers.

\subsection{Side discovery: reflection and Snell's laws}
At the plane $z=0$, the boundary condition must hold for every $x$ and time:
\begin{equation}
E_i\ee^{\ii k_{x,i}x-\ii\omega_i t}
+E_r\ee^{\ii k_{x,r}x-\ii\omega_r t}
=E_t\ee^{\ii k_{x,t}x-\ii\omega_t t}.
\end{equation}
For a stationary, uniform, linear planar interface, a single incident Fourier component cannot produce a different surface frequency or transverse wave number. Hence
\begin{equation}
\omega_i=\omega_r=\omega_t,\qquad
k_{x,i}=k_{x,r}=k_{x,t},\qquad k_{y,i}=k_{y,r}=k_{y,t}.
\end{equation}
The reflected longitudinal component changes sign in medium 0, giving the familiar equality of reflection and incidence angles. Since $k_x=k_0\bn\sin\theta$, the transmitted component obeys
\begin{equation}
\boxed{\bn_0\sin\theta_0=\bn_1\sin\theta_1.}
\end{equation}
This is Snell's law recovered directly from boundary phase matching: a side discovery of a previously known geometrical result.

% For absorbing materials, a complex angle is an algebraic shorthand. It is safer numerically to compute
% \begin{equation}
% k_{z,j}=\sqrt{k_0^2\bn_j^2-k_\parallel^2}.
% \label{eq:kz}
% \end{equation}
% Choose the passive outgoing branch: $\Ima k_{z,j}\geq0$, and the forward-energy branch in the lossless propagating limit. An evanescent transmitted field has positive imaginary $k_z$ and decays. A grating can change tangential momentum by reciprocal-lattice vectors; the uniform-interface argument does not forbid diffraction at a patterned mask.

\subsection{Why the reflected magnetic field has a minus sign}
For a plane wave, $\nabla\to\ii\mathbf k$ in Faraday's law:
\begin{equation}
\boxed{\mathbf H=\frac{1}{\mu_0\omega}\mathbf k\times\mathbf E.}
\label{eq:kcrossE}
\end{equation}
At normal incidence, $\mathbf k\parallel+\widehat z$ and $\mathbf E\parallel\widehat x$ give $\mathbf H\parallel+\widehat y$. Reversing $\mathbf k$ while retaining the fixed electric axis reverses $\mathbf H$. Define $Y_j$ as the forward-wave ratio of the suitably oriented tangential magnetic field to tangential electric field. Then
\begin{equation}
H_i=Y_0E_i,\qquad H_r=-Y_0E_r,\qquad H_t=Y_1E_t.
\end{equation}
Magnetic continuity therefore gives
\begin{equation}
Y_0E_i-Y_0E_r=Y_1E_t
\quad\Longrightarrow\quad\boxed{Y_0(1-r)=Y_1t.}
\label{eq:Hbc}
\end{equation}
The minus sign expresses the reflected propagation direction; it is not a phenomenological absorption term.

\subsection{Solving for reflection and transmission}
Eliminate $t$ using \eqref{eq:Ebc}:
\begin{align}
Y_0(1-r)&=Y_1(1+r),\\
Y_0-Y_1&=(Y_0+Y_1)r.
\end{align}
Thus
\begin{equation}
\boxed{r_{01}=\frac{Y_0-Y_1}{Y_0+Y_1},\qquad
 t_{01}=1+r_{01}=\frac{2Y_0}{Y_0+Y_1}.}
\label{eq:rt}
\end{equation}
If $Y_0=Y_1$, then $r=0,t=1$: the incident wave already satisfies both boundary conditions. A mismatch generally requires a reflected wave to make both fields match.

\subsection{Secondary identities for the later film sum}
Interchanging the incident and transmitted media, while keeping the same tangential-field convention,
\begin{equation}
r_{10}=\frac{Y_1-Y_0}{Y_0+Y_1}=-r_{01},\qquad
 t_{10}=\frac{2Y_1}{Y_0+Y_1}.
\end{equation}
The product and square satisfy
\begin{align}
t_{01}t_{10}&=\frac{4Y_0Y_1}{(Y_0+Y_1)^2},\\
1-r_{01}^2&=\frac{(Y_0+Y_1)^2-(Y_0-Y_1)^2}{(Y_0+Y_1)^2}
=\frac{4Y_0Y_1}{(Y_0+Y_1)^2}.
\end{align}
Therefore
\begin{equation}
\boxed{r_{10}=-r_{01},\qquad t_{01}t_{10}=1-r_{01}^2.}
\label{eq:identities}
\end{equation}
These are algebraic identities for complex amplitudes. The square is $r^2$, \emph{not} $|r|^2$. They remain valid for complex admittances within the stated interface model.

\subsection{Field coefficients versus power coefficients}
For a forward wave the normal time-averaged energy flux is
$S_z=\tfrac12\Rea(Y)|E_\tau|^2$. With a lossless incident medium and flux evaluated at the interface,
\begin{equation}
\boxed{\mathcal R=|r|^2,\qquad
\mathcal T=\frac{\Rea Y_1}{Y_0}|t|^2.}
\end{equation}
For real positive $Y_0,Y_1$, substitution into \eqref{eq:rt} verifies
$\mathcal R+\mathcal T=[(Y_0-Y_1)^2+4Y_0Y_1]/(Y_0+Y_1)^2=1$.
If medium 1 absorbs, $\mathcal T$ here is the flux entering it, which then decays with depth. In an absorbing incident medium, incident/reflected interference also contributes to net normal flux; the naive identification of a separately normalized $1-\mathcal R$ with net input needs care.

\section{Polarization-specific admittances and Fresnel formulas}
Let the incidence plane be $xz$ and the interface normal be $z$.
The \textbf{vacuum wave impedance} is
\begin{equation}
\boxed{Z_{\rm vac}=\sqrt{\frac{\mu_0}{\epsilon_0}}=\mu_0c}
% =\frac{1}{\epsilon_0c}\simeq376.73\,\Omega.}
\end{equation}
Here $\mu_0$ and $\epsilon_0$ are the vacuum permeability and permittivity, and $c=1/\sqrt{\mu_0\epsilon_0}$ is the speed of light in vacuum. For a travelling plane wave in vacuum, $Z_{\rm vac}$ is the ratio of electric- to magnetic-field amplitudes, $|\mathbf E|/|\mathbf H|$; its reciprocal is the vacuum wave admittance. The polarization-specific $Y$ below instead relates the oriented tangential magnetic field to the tangential electric field, so it also contains the incidence-angle factors.
For nonmagnetic media, $\epsilon=\epsilon_0\bn^2$.

\subsection{s (TE) polarization: derive its admittance}
Here $\mathbf E=(0,E_y,0)$ and $\mathbf k=(k_x,0,k_z)$. From \eqref{eq:kcrossE},
\begin{equation}
\mathbf H=\frac{1}{\mu_0\omega}(-k_zE_y,0,k_xE_y).
\end{equation}
Choose $-H_x$ as the oriented tangential magnetic component of a forward wave. Then
\begin{equation}
\boxed{Y_s=\frac{-H_x}{E_y}=\frac{k_z}{\mu_0\omega}
=\frac{\bn\cos\theta}{Z_{\rm vac}}.}
\end{equation}
A backward wave has $k_z\to-k_z$, giving the required opposite sign.

\subsection{p (TM) polarization: derive its admittance}
Now $\mathbf E=(E_x,0,E_z)$ and $\mathbf H=(0,H_y,0)$. In a homogeneous source-free isotropic medium, $\mathbf k\cdot\mathbf D=0$ gives
\begin{equation}
k_xE_x+k_zE_z=0\quad\Longrightarrow\quad E_z=-\frac{k_x}{k_z}E_x.
\end{equation}
For complex $\mathbf k$, this is a bilinear algebraic transversality condition, not ordinary Euclidean orthogonality to a real ray direction. Substitution gives
\begin{align}
H_y&=\frac{k_zE_x-k_xE_z}{\mu_0\omega}
=\frac{k_z^2+k_x^2}{\mu_0\omega k_z}E_x
=\frac{k_0^2\bn^2}{\mu_0\omega k_z}E_x,\\
Y_p&=\frac{H_y}{E_x}=\frac{\omega\epsilon}{k_z}
=\boxed{\frac{\bn}{Z_{\rm vac}\cos\theta}}.
\end{align}
Here the amplitude is $E_x$, the tangential electric component. Using the full $p$-field amplitude instead changes coefficient formulas.

\subsection{The polarization-specific coefficients}
Insert each admittance in \eqref{eq:rt}. The common $1/Z_{\rm vac}$ cancels:
\begin{align}
r_s&=\frac{\bn_0\cos\theta_0-\bn_1\cos\theta_1}
{\bn_0\cos\theta_0+\bn_1\cos\theta_1},&
t_s&=\frac{2\bn_0\cos\theta_0}
{\bn_0\cos\theta_0+\bn_1\cos\theta_1},\\
r_{p,\tau}&=\frac{\bn_0\cos\theta_1-\bn_1\cos\theta_0}
{\bn_0\cos\theta_1+\bn_1\cos\theta_0},&
t_{p,\tau}&=\frac{2\bn_0\cos\theta_1}
{\bn_0\cos\theta_1+\bn_1\cos\theta_0}.
\end{align}
The $p$ expressions follow by multiplying numerator and denominator of the admittance ratios by $\cos\theta_0\cos\theta_1$. At normal incidence both tangential-field polarizations give
\begin{equation}
r=\frac{\bn_0-\bn_1}{\bn_0+\bn_1},\qquad
 t=\frac{2\bn_0}{\bn_0+\bn_1}.
\end{equation}
For real propagating waves, the commonly used reflected $p$ basis may have a reversed tangential axis. In that convention $r_p^{\rm conventional}=-r_{p,\tau}$ and
$t_p^{\rm full}=(\cos\theta_0/\cos\theta_1)t_{p,\tau}$. This sign difference is a basis choice; it must not be changed halfway through a multilayer sum.

\textbf{A check: Brewster incidence.} For lossless positive indices, $Y_{p,0}=Y_{p,1}$ means $n_0\cos\theta_1=n_1\cos\theta_0$. Combining with Snell's law gives $\tan\theta_B=n_1/n_0$ (for distinct indices). Absorbing media generally have a reflectance minimum rather than an exactly zero real-angle reflection.
\clearpage
